% Options for packages loaded elsewhere
\PassOptionsToPackage{unicode}{hyperref}
\PassOptionsToPackage{hyphens}{url}
%
\documentclass[
]{article}
\usepackage{amsmath,amssymb}
\usepackage{iftex}
\ifPDFTeX
  \usepackage[T1]{fontenc}
  \usepackage[utf8]{inputenc}
  \usepackage{textcomp} % provide euro and other symbols
\else % if luatex or xetex
  \usepackage{unicode-math} % this also loads fontspec
  \defaultfontfeatures{Scale=MatchLowercase}
  \defaultfontfeatures[\rmfamily]{Ligatures=TeX,Scale=1}
\fi
\usepackage{lmodern}
\ifPDFTeX\else
  % xetex/luatex font selection
\fi
% Use upquote if available, for straight quotes in verbatim environments
\IfFileExists{upquote.sty}{\usepackage{upquote}}{}
\IfFileExists{microtype.sty}{% use microtype if available
  \usepackage[]{microtype}
  \UseMicrotypeSet[protrusion]{basicmath} % disable protrusion for tt fonts
}{}
\makeatletter
\@ifundefined{KOMAClassName}{% if non-KOMA class
  \IfFileExists{parskip.sty}{%
    \usepackage{parskip}
  }{% else
    \setlength{\parindent}{0pt}
    \setlength{\parskip}{6pt plus 2pt minus 1pt}}
}{% if KOMA class
  \KOMAoptions{parskip=half}}
\makeatother
\usepackage{xcolor}
\usepackage[margin=1in]{geometry}
\usepackage{graphicx}
\makeatletter
\def\maxwidth{\ifdim\Gin@nat@width>\linewidth\linewidth\else\Gin@nat@width\fi}
\def\maxheight{\ifdim\Gin@nat@height>\textheight\textheight\else\Gin@nat@height\fi}
\makeatother
% Scale images if necessary, so that they will not overflow the page
% margins by default, and it is still possible to overwrite the defaults
% using explicit options in \includegraphics[width, height, ...]{}
\setkeys{Gin}{width=\maxwidth,height=\maxheight,keepaspectratio}
% Set default figure placement to htbp
\makeatletter
\def\fps@figure{htbp}
\makeatother
\ifLuaTeX
  \usepackage{luacolor}
  \usepackage[soul]{lua-ul}
\else
  \usepackage{soul}
\fi
\setlength{\emergencystretch}{3em} % prevent overfull lines
\providecommand{\tightlist}{%
  \setlength{\itemsep}{0pt}\setlength{\parskip}{0pt}}
\setcounter{secnumdepth}{5}
\ifLuaTeX
  \usepackage{selnolig}  % disable illegal ligatures
\fi
\usepackage{bookmark}
\IfFileExists{xurl.sty}{\usepackage{xurl}}{} % add URL line breaks if available
\urlstyle{same}
\hypersetup{
  pdftitle={Problem Set 1 - Predicción del ingreso},
  pdfauthor={Maria Camila Caraballo, Javier Amaya, Nicolas Moreno, Mateo Isaza},
  hidelinks,
  pdfcreator={LaTeX via pandoc}}

\title{Problem Set 1 - Predicción del ingreso}
\author{Maria Camila Caraballo, Javier Amaya, Nicolas Moreno, Mateo
Isaza}
\date{2025-02-01}

\begin{document}
\maketitle

{
\setcounter{tocdepth}{2}
\tableofcontents
}
\newpage{}

\section{4. Brecha salarial por género (Javier /
Mateo)}\label{brecha-salarial-por-guxe9nero-javier-mateo}

\subsection{Rubrica}\label{rubrica}

Se presenta la estimación de la brecha salarial de género incondicional
y condicional en una misma tabla, con las estimaciones lado a lado. La
tabla no es un pantallazo, sino que está debidamente formateada, con
título y notas que permitan entenderla sin recurrir al documento. Por
ejemplo, listando las variables control utilizadas. Se justifica la
inclusión de los diferentes controles.

\subsection{Descripción taller}\label{descripciuxf3n-taller}

La brecha salarial de género. Los formuladores de políticas han estado
preocupados durante mucho tiempo por la brecha salarial de género, y
este será nuestro enfoque en esta subsección.

\subsubsection{\texorpdfstring{\ul{a. Comience estimando y discutiendo
la brecha salarial
incondicional:}}{a. Comience estimando y discutiendo la brecha salarial incondicional:}}\label{a.-comience-estimando-y-discutiendo-la-brecha-salarial-incondicional}

\[
log(w) = \beta_1 + \beta_2Female + u
\]

donde \emph{Female} es un indicador que toma el valor de uno si el
individuo en la muestra se identifica como mujer.

\% Table created by stargazer v.5.2.3 by Marek Hlavac, Social Policy
Institute. E-mail: marek.hlavac at gmail.com \% Date and time: sáb.,
mar. 01, 2025 - 3:00:40 p.~m. \% Requires LaTeX packages: dcolumn

\begin{table}[!htbp] \centering 
  \caption{Tabla XXX. Resultados 'unconditional model' para brecha salarial} 
  \label{} 
\begin{tabular}{@{\extracolsep{5pt}}lD{.}{.}{-3} } 
\\[-1.8ex]\hline 
\hline \\[-1.8ex] 
 & \multicolumn{1}{c}{\textit{Dependent variable:}} \\ 
\cline{2-2} 
\\[-1.8ex] & \multicolumn{1}{c}{log(ingreso)} \\ 
\hline \\[-1.8ex] 
 Female & -0.062 \\ 
  & (0.013) \\ 
 \hline \\[-1.8ex] 
Observations & \multicolumn{1}{c}{16,542} \\ 
R$^{2}$ & \multicolumn{1}{c}{0.001} \\ 
Adjusted R$^{2}$ & \multicolumn{1}{c}{0.001} \\ 
\hline 
\hline \\[-1.8ex] 
\textit{Note:}  & \multicolumn{1}{l}{Standard errors in parentheses.} \\ 
\end{tabular} 
\end{table}

\textbf{respuesta:}

Utilizando como referencia el ``unconditional model'' podemos discutir
sobre los siguientes puntos:

\begin{itemize}
\item
  Para el caso del intercepto, podemos sacarle el exponente para obtener
  el salario promedio de los hombres que llega a \$ 6314.9 por hora.
\item
  En cuanto al coeficiente de \textbf{\emph{female}} en el que se puede
  observar que dicho coeficiente es estadisticamente diferente de cero
  (\emph{p value= 0.00001}), podemos interpretar que en promedio el
  salario por hora de las mujeres es un 6.2\% más bajo que el salario de
  los hombres.
\item
  En cuanto al \(R^2\) ajustado se puede observar que es del 1\% lo que
  quiere decir que el modelo utilizando solo variable de \emph{female}
  explica el 1\% de la variabilidad del logaritmo del ingreso por hora.
  De alguna manera, esta es esperable porque solo tenemos una variable
  predictora y esto deja por fuera otras variables que pueden ser
  relevantes como la educación, la edad, la experiencia, entre otras.
\end{itemize}

\subsubsection{\texorpdfstring{\ul{b. Un eslogan común es ``igual
salario por igual trabajo''. Una forma de interpretar esto es que, para
empleados con características laborales y personales similares, no
debería existir una brecha salarial de género. Para analizar esto,
estima la brecha salarial condicional incorporando variables de control
como características del trabajador y del puesto de
trabajo.}}{b. Un eslogan común es ``igual salario por igual trabajo''. Una forma de interpretar esto es que, para empleados con características laborales y personales similares, no debería existir una brecha salarial de género. Para analizar esto, estima la brecha salarial condicional incorporando variables de control como características del trabajador y del puesto de trabajo.}}\label{b.-un-eslogan-comuxfan-es-igual-salario-por-igual-trabajo.-una-forma-de-interpretar-esto-es-que-para-empleados-con-caracteruxedsticas-laborales-y-personales-similares-no-deberuxeda-existir-una-brecha-salarial-de-guxe9nero.-para-analizar-esto-estima-la-brecha-salarial-condicional-incorporando-variables-de-control-como-caracteruxedsticas-del-trabajador-y-del-puesto-de-trabajo.}

\textbf{Respuesta:} para llevar a cabo esta estimación de un modelo
condicional se consideraron dos grandes grupos de variables control que
pueden afectar la estimación de la relación de asociación entre el
salario por hora y el género. El primer grupo son las variables
relacionadas con las características del trabajador como el nivel de
educación alcanzado (a mayores niveles de estudio se espera una mayor
remuneración y el nivel de estudio máximo alcanzado puede ser diferente
en mujeres vs hombres), la edad (algunas edades están relacionadas con
mayores ingresos como personas entre 30 y 50 años y esto puede diferir
entre hombres y mujeres por los patrones de
embarazo-lactancia-maternidad) y la antiguedad del trabajador en la
empresa (se espera que trabajadores con mayor antiguedad tengan un mayor
ingreso). En el segundo grupo de variables se pueden encontrar variables
como el tipo de oficio (algunas ocupaciones de predominio de hombres
como la ingenieria tienden a tener mayotes salarios); informalidad (en
algunos países las mujeres tienden a estar sobre representadas en el
grupo de trabajadores informales, trabajos que usualmente reciben menos
ingresos) y el tamaño de la empresa (puede ocurrir que en empresas
pequeñas y medianas haya una mayor proporción de mujeres y en empresas
grandes haya mayor proporción de hombres, siendo estas últimas lugares
donde se espera mejores salarios).

\begin{enumerate}
\def\labelenumi{\roman{enumi}.}
\tightlist
\item
  Primero estimamos la brecha salarial condicional a los controles
  explicados previamente usando el método de FWL:
\end{enumerate}

\% Table created by stargazer v.5.2.3 by Marek Hlavac, Social Policy
Institute. E-mail: marek.hlavac at gmail.com \% Date and time: sáb.,
mar. 01, 2025 - 3:00:47 p.~m. \% Requires LaTeX packages: dcolumn

\begin{table}[!htbp] \centering 
  \caption{Tabla XXX. Resultados 'unconditional model' para brecha salarial} 
  \label{} 
\begin{tabular}{@{\extracolsep{5pt}}lD{.}{.}{-4} D{.}{.}{-4} D{.}{.}{-4} } 
\\[-1.8ex]\hline 
\hline \\[-1.8ex] 
 & \multicolumn{3}{c}{Logaritmo del ingreso} \\ 
\cline{2-4} 
 & \multicolumn{1}{c}{Simple} & \multicolumn{1}{c}{Múltiple} & \multicolumn{1}{c}{FWL} \\ 
\\[-1.8ex] & \multicolumn{1}{c}{(1)} & \multicolumn{1}{c}{(2)} & \multicolumn{1}{c}{(3)}\\ 
\hline \\[-1.8ex] 
 Female & -0.0615 & -0.1076 & -0.1076 \\ 
  & (0.0133) & (0.0121) & (0.0120) \\ 
 \hline \\[-1.8ex] 
Observations & \multicolumn{1}{c}{16,542} & \multicolumn{1}{c}{16,542} & \multicolumn{1}{c}{16,542} \\ 
R$^{2}$ & \multicolumn{1}{c}{0.0013} & \multicolumn{1}{c}{0.4500} & \multicolumn{1}{c}{0.0048} \\ 
Adjusted R$^{2}$ & \multicolumn{1}{c}{0.0012} & \multicolumn{1}{c}{0.4378} & \multicolumn{1}{c}{0.0048} \\ 
\hline 
\hline \\[-1.8ex] 
\textit{Note:}  & \multicolumn{3}{l}{Standard errors in parentheses.} \\ 
\end{tabular} 
\end{table}

\end{document}
